\subsubsection{BDSM 后的情绪跌落（Sub Drop / Top Drop）：机制、识别与照护}

本书已多次提到 aftercare（事后关怀）的做法——拥抱、饮水、交谈、复盘（参见 SM 健康与安全相关小节及书末词汇表）。但还有一个常被漏掉的现象：\textbf{情绪跌落（drop）}——强烈的 BDSM 活动结束后数小时到两三天内，参与者陷入无来由的低落、疲惫、空虚、焦虑或自我怀疑。受方（sub drop）与施方（top drop）都会经历。

\textbf{机制}：高强度场景中，内啡肽、多巴胺、催产素与肾上腺素处于高峰——这是"飞行状态"的来源；活动结束后，它们回落到基线水平，甚至短暂跌破基线（类似长跑冲线后的低谷、演出谢幕后的失落）。这不是"玩出问题了"，而是神经化学的正常账单。

\textbf{两个容易误判的特征}：
\begin{itemize}
\item \textbf{延迟出现}：当晚一切正常，第二、三天突然情绪跌入谷底，伴随莫名的羞耻感或"我是不是有问题"的自我怀疑——很多人因此错误归因为"关系出问题了"或"那次游戏越界了"，实际只是跌落期；
\item \textbf{施方同样会跌落}：top drop 常被忽视——施虐方同样经历神经化学回落，还可能叠加"我刚才是不是太过分了"的道德反刍。
\end{itemize}

\textbf{照护安排}（把 drop 预案提升到与安全词同级的地位）：
\begin{itemize}
\item 当晚 aftercare 做足：保暖、进食、补水、肯定的言语——"刚才的场景我们都在状态里"比沉默刷手机更能托住情绪；
\item 预约复盘：约定 48-72 小时内再做一次 check-in（一条消息即可），专门用于核对跌落期情绪——被预先告知"可能会低落"的大脑，真正低落时的恐慌会小得多；
\item 生理支持：跌落期多晒太阳、轻度运动、规律进食与睡眠，帮助神经递质回升；\textbf{跌落期不做重大关系决定，不解读关系走向}；
\item 求助边界：低落持续两周以上、影响睡眠饮食与日常功能时，按抑郁的常规筛查就医处理，不再归因于玩法。
\end{itemize}

\noindent\textit{【编者注】}drop 的神经化学解释目前主要来自一般应激科学、运动科学的外推与社群经验总结，针对性的对照研究仍然有限；但该现象的普遍性已被社群与临床观察广泛报告，上述照护建议与已知机制不冲突，可安全采纳。
